\documentclass[oneside,senior,etd]{BYUPhysForDegree}

\usepackage{fontspec}
\setmainfont{Times New Roman}
\usepackage{rotating} 
\usepackage{dsfont}

\usepackage[russian]{babel}
\usepackage{amsfonts} % Пакеты для математических символов и теорем
\usepackage{amstext}
\usepackage{amssymb}
\usepackage{amsthm}
\usepackage{graphicx} % Пакеты для вставки графики
\usepackage{subfig}
\usepackage{color}
\usepackage{subcaption}
\usepackage[unicode]{hyperref}
\usepackage[nottoc]{tocbibind}
\usepackage[ruled,vlined]{algorithm2e}
\SetAlFnt{\small} % размер шрифта для алгоритма
\SetAlCapFnt{\small} % размер шрифта для заголовка
\SetAlCapNameFnt{\small} % размер шрифта для названия
\usepackage{verbatim} % Для вставок заранее подготовленного текста в режиме as-is
\usepackage{listings}

\usepackage{commath}
\newcommand\Tau{\mathcal{T}}
\newcommand{\R}{\mathbb{R}}
\usepackage{color}
\usepackage[colorinlistoftodos, prependcaption]{todonotes}
\usepackage{multirow}
\usepackage{indentfirst}
\newcommand*{\MyIndent}{\hspace*{0.2cm}}%

\usepackage{pdfpages}

\Chair{Кафедра интеллектуальных информационных технологий}
\Lab{~}
\Year{2026}
  \Month{Июнь}
  \City{Москва}
  \AuthorText{Автор}
  \Author{Трифонов Андрей Рафисович}
  \AuthorEng{}
  \AcadGroup{622}
  % Если вы не Пупкин Василий Иванович, то скопируйте себе этот проект через меню и уже потом правьте
  \TitleTop{Ускорение алгоритмов растеризации}
  %\TitleMiddle{}
  \TitleBottom{поверхностей, заданных неявными функциями,\\на графических процессорах}
  %\TitleMiddle{}
  \TitleTopEng{TODO}
  \TitleBottomEng{} % leave empty if you don't need it  
  %\docname{Курсовая работа}
  \docname{Магистерская диссертация}
  %\docname{Магистерская диссертация}
  \Advisor{Фролов Владимир Александрович}  
  \AdvisorDegree{к.ф.-м.н.}
  %\Consultant{}
  %\ConsultantDegree{}
  


%%%% DON'T change this. It is here because .sty does not support cyrillic cp properly %%%%
\University{Московский государственный университет имени М.В.Ломоносова}
\Faculty{Факультет вычислительной математики и кибернетики}
\GrText{гр.}
\AdvisorText{Научный руководитель}
\ConsultantText{Научный консультант}
\AbstractText{Аннотация}

\begin{document}
\makepreliminarypages
\fixmargins

\oneandhalfspace

\tableofcontents

\newpage
\section{Введение}
\label{sec:Chapter0} \index{Chapter0}

\subsection{Введение в 3D-графику и её задачи}
Современный мир невозможно представить без 3D-графики, которая прочно вошла во все сферы нашей жизни и стала одним из ключевых инструментов визуализации, моделирования и взаимодействия с данными. От интерактивных видеоигр и кинематографических спецэффектов до сложных научных симуляций, медицинских визуализаций, систем автоматизированного проектирования (САПР) и стремительно развивающихся технологий виртуальной и дополненной реальности (VR/AR),— трехмерные изображения играют центральную роль в представлении сложных объектов, процессов и сред.

С развитием технологий и ростом вычислительных мощностей графических процессоров (GPU) постоянно возрастают требования к качеству, детализации и интерактивности 3D-изображений. Пользователи ожидают фотореалистичной графики, мгновенного отклика и возможности взаимодействия со сложными трехмерными сценами в реальном времени. Это создает постоянный вызов для разработчиков алгоритмов и систем 3D-графики, требуя поиска новых, более эффективных и гибких подходов к моделированию и, что особенно важно, к рендерингу. Особую актуальность приобретают методы, способные обеспечить высокую производительность при работе с высокодетализированными и динамически изменяющимися сценами, стремясь при этом к максимальному реализму.

\subsection{Классические подходы к представлению 3D-моделей}

На протяжении многих десятилетий доминирующим и наиболее распространенным способом представления трехмерных объектов в компьютерной графике остается использование полигональных моделей (также известных как сетки, или mesh-модели). В этом подходе поверхность объекта аппроксимируется набором плоских геометрических примитивов – обычно треугольников или четырехугольников (квадов). Каждая вершина имеет координаты в трехмерном пространстве, а также может содержать дополнительную информацию, такую как нормали к поверхности аппроксимируемого объекта, координаты текстур и цвет.

Преимущества полигональных моделей:

\begin{itemize}
    \item 
    \textbf{Простота математического описания:} Треугольник – это простейший плоский многоугольник, математические операции с которым (интерполяция, пересечение луча и т.д.) относительно просты и хорошо изучены.
    \item 
    \textbf{Хорошая поддержка аппаратным обеспечением:} Современные графические процессоры (GPU) изначально спроектированы и оптимизированы для эффективной обработки и растеризации полигональных моделей. Это позволяет достигать высокой производительности рендеринга.
    \item 
    \textbf{Развитые инструменты и стандарты:} Существует огромное количество программного обеспечения для 3D-моделирования (например, Blender, Autodesk Maya, 3ds Max) и межплатформенных форматов (например, OBJ, FBX, GLTF), которые активно используются и постоянно развиваются индустрией.
    \item 
    \textbf{Простота визуализации:} Прямая растеризация полигонов через стандартный графический конвейер (pipeline) обеспечивает высокую скорость вывода изображения на экран.
\end{itemize}

Однако, несмотря на свои преимущества, полигональные модели имеют ряд существенных ограничений, особенно при работе со сложными, органическими или динамически изменяющимися поверхностями:

\begin{itemize}
    \item 
    \textbf{Топологические ограничения:} Слишком плотная или, наоборот, разреженная полигональная сетка может создавать проблемы при моделировании и деформации. Создание плавных, округлых форм требует большого количества полигонов, что увеличивает объем данных.
    \item 
    \textbf{Проблема уровня детализации (Level of Detail, LOD):} Для эффективного рендеринга объектов на разных расстояниях от камеры (или с разной степенью важности) необходимо иметь несколько версий одной и той же модели с различным числом полигонов. Управление этими версиями и переключение между ними требует дополнительных усилий и усложняет конвейер разработки.
    \item 
    \textbf{Сложность описания органических форм и динамических изменений:} Точное представление объектов со сложной, изменяющейся топологией (например, жидкости, облака, мягкие тела, волосы) с помощью фиксированной полигональной сетки является крайне сложной или невозможной задачей без постоянного перестроения сетки.
    \item 
    \textbf{Булевы операции:} Выполнение таких операций, как объединение, пересечение или вычитание полигональных моделей, является нетривиальной задачей и часто приводит к созданию некачественной (non-manifold, с самопересечениями) или избыточной геометрии, требующей ручной доработки.
    \item 
    \textbf{Ограничения по памяти:} Для достижения высокой детализации требуется экспоненциально увеличивать количество полигонов, что приводит к большим объемам данных и потенциальным проблемам с хранением, передачей и обработкой.
\end{itemize}

Эти ограничения обусловили поиск альтернативных способов представления 3D-моделей, которые могли бы обеспечить большую гибкость и эффективность, особенно в условиях возрастающих требований к сложности и динамичности трехмерных сцен.

\subsection{Неявные поверхности: функции расстояния}
Ограничения полигонального представления, особенно в контексте моделирования сложных органических форм, выполнения булевых операций и потребности в адаптивном уровне детализации, стимулировали развитие альтернативных методов представления трехмерных объектов. Одним из таких направлений стало использование неявных поверхностей (Implicit Surfaces).

В отличие от явного (explicit) представления, где поверхность описывается параметрически в виде набора точек, линий или полигонов (например, полигональные сетки, NURBS-кривые \cite{piegl2012nurbs}), неявная поверхность определяется как множество точек (x,y,z), для которых некоторая функция F(x,y,z) обращается в ноль (то есть F(x,y,z)=0).

Особое значение среди неявных поверхностей получили функции расстояния (Distance Functions), а в частности, знаковые функции расстояния (Signed Distance Fields, SDFs) \cite{Osher2003}. SDF представляет собой функцию D(x,y,z), которая для любой точки в пространстве возвращает кратчайшее евклидово расстояние до ближайшей точки поверхности объекта. При этом знак функции указывает, находится ли точка внутри объекта (отрицательное значение) или снаружи (положительное значение). Нулевой уровень этой функции (D(x,y,z)=0) точно определяет поверхность объекта.

Ключевые преимущества SDFs:

\begin{itemize}
    \item 
    \textbf{Компактное представление:} Сложные топологии и детализированные поверхности могут быть описаны относительно простой аналитической функцией или дискретным набором значений, хранимых в иерархических структурах данных (например, окдеревьях, воксельных сетках), вместо огромного числа полигонов. Это обеспечивает высокую экономию памяти для хранения.
    \item 
    \textbf{Гибкость в создании и модификации:} SDFs позволяют легко выполнять булевы операции (объединение, пересечение, вычитание) над сложными формами простыми математическими операциями. Это значительно упрощает процедурное моделирование и создание сложных объектов.
    \item 
    \textbf{"Бесконечная" детализация:} Поскольку поверхность определяется функцией, а не дискретным набором данных, возможно получение любого уровня детализации "на лету" без потери информации. Для увеличения детализации достаточно просто увеличить количество "запросов" к функции расстояния.
    \item 
    \textbf{Простота получения нормалей:} В любой точке поверхности вектор нормали может быть легко вычислен как градиент функции расстояния.
    \item 
    \textbf{Наличие информации о расстоянии до поверхности:} Знать расстояние до поверхности полезно для определения столкновений, расчета толщины объекта, сглаживания краев и других применений.
    \item 
    \textbf{Использование в различных областях:} SDFs находят широкое применение в процедурной генерации контента (пейзажи, персонажи, текстуры), симуляции жидкостей и мягких тел, 3D-печати (для проверки корректности геометрии), быстрого прототипирования, а также для хранения геометрии в задачах машинного обучения.
\end{itemize}

Несмотря на свои значительные преимущества, использование SDFs сопряжено с рядом фундаментальных проблем, особенно в контексте интерактивной визуализации:

\begin{itemize}
    \item 
    \textbf{Несовместимость с традиционным полигональным графическим конвейером:} Стандартные GPU-конвейеры разработаны для рендеринга полигональных моделей. Прямой рендеринг SDFs требует специализированных алгоритмов, таких как трассировка лучей (Ray Marching), которые могут быть вычислительно затратными.
    \item 
    \textbf{Сложность эффективной конвертации в полигональную сетку в реальном времени:} Хотя методы вроде Marching Cubes позволяют преобразовать SDF в полигоны, они часто требуют дискретизации пространства (вокселизации) и сопряжены с проблемами потери детализации или генерации избыточного количества геометрии. Создание адаптивной и высококачественной полигональной сетки из SDF в реальном времени остается сложной задачей.
\end{itemize}

Эти вызовы подчеркивают необходимость разработки новых, высокопроизводительных и адаптивных методов рендеринга SDFs, способных эффективно использовать преимущества данного представления и соответствовать требованиям современных графических приложений.

\subsection{Актуальность диссертации}
В последние годы, с бурным развитием графических процессоров и появлением новых аппаратных архитектур, таких как Mesh Shaders \cite{NVIDIA_MeshShaders} в современных GPU, открываются принципиально новые возможности для переосмысления процессов рендеринга неявных поверхностей. Mesh Shaders позволяют программировать конвейер генерации геометрии непосредственно на GPU, обеспечивая беспрецедентную гибкость и управляемость на ранних стадиях графического конвейера. Это создает уникальную возможность для реализации высокоэффективных алгоритмов динамической триангуляции неявных поверхностей, адаптированных под современное аппаратное обеспечение.

Таким образом, актуальность данной магистерской диссертации определяется острой потребностью в разработке и оптимизации методов рендеринга 3D-моделей, заданных функциями расстояния, которые могли бы:

\begin{itemize}
    \item 
     Эффективно использовать современные аппаратные возможности GPU, в частности, Mesh Shaders, для максимально производительной и адаптивной генерации геометрии "на лету".
    \item 
     Снизить вычислительные затраты и объем генерируемой геометрии за счет адаптивной триангуляции видимых участков поверхности.
\end{itemize}

Разработка такого метода, базирующегося на динамической триангуляции и специализированных GPU-конвейерах, позволит преодолеть ограничения существующих подходов и сделает интерактивную визуализацию сложных SDF-моделей более доступной и производительной, открывая новые горизонты для их применения в областях, где сейчас доминирует полигональное представление.

\section{Постановка задачи}
\subsection{Цель}
Разработать быстрый метод рендеринга неявных представлений трёхмерных моделей на основе SDF.
\subsection{Задачи}
\begin{enumerate}
\item Провести обзор существующих неявных представлений и выбрать базовое представление и базовый метод рендеринга.
\item Разработать алгоритм адаптивной растеризации.
\item Ускорить разработанный алгоритм с использованием аппаратных возможностей современных графических процессоров.
\item Провести его экспериментальную оценку и сравнение с аналогами.
\end{enumerate}

\subsection{Формальная постановка задачи}
\subsubsection{Входные данные}
\begin{itemize}
    \item \textbf{Функция расстояния (SDF) $D(\mathbf{p})$:} Знаковая функция расстояния \cite{Osher2003}, где $\mathbf{p} = (x, y, z) \in \mathbb{R}^3$.
    $D(\mathbf{p}) \in \mathbb{R}$ обозначает кратчайшее евклидово расстояние от $\mathbf{p}$ до поверхности объекта $S = \{\mathbf{p} \in \mathbb{R}^3 \mid D(\mathbf{p}) = 0\}$.
    \item \textbf{Параметры камеры $C$:} Набор параметров $C = (Pos, Dir, Up, Fov, Aspect)$, где $Pos \in \mathbb{R}^3$ (позиция), $Dir \in \mathbb{R}^3$ (направление взгляда), $Up \in \mathbb{R}^3$ (вектор "вверх"), $Fov \in (0, \pi)$ (угол обзора), $Aspect \in \mathbb{R}^+$ (соотношение сторон).
    \item \textbf{Параметры рендеринга $P$:} Набор управляющих параметров $P = \{T_{detail}, T_{adaptive}, \dots\}$, где $T_{detail}$ - желаемый уровень детализации, $T_{adaptive}$ - параметры адаптивности.
\end{itemize}

\subsubsection{Выходные данные}
\begin{itemize}
    \item \textbf{Изображение $I$:} Двумерное цветное (трёхканальное) изображение $I \in \mathbb{R}^{W_{scr} \times H_{scr} \times 3}$, где $W_{scr}, H_{scr}$ - ширина и высота экрана в пикселях.
\end{itemize}

\subsubsection{Задача}

Найти отображение $\mathcal{G}: (\mathcal{D}_{SDF}, \mathcal{C}_{params}, \mathcal{P}_{params}) \rightarrow \mathcal{I}_{output}$ такое, что на основе входных данных (функции расстояния $D(\mathbf{p})$, параметров камеры $C$, параметров рендеринга $P$) будет сформировано изображение $I$.
Основной процесс $\mathcal{G}$ включает:
\begin{enumerate}
    \item \textbf{Динамическую триангуляцию:} Преобразование $D(\mathbf{p})$ в адаптивную полигональную сетку $\Delta = \{\text{tri}_j\}_{j=1}^M$, где $\text{tri}_j$ - треугольник, аппроксимирующий поверхность $S$;
    \item \textbf{Рендеринг:} Визуализация $\Delta$ в изображение $I$ через графический конвейер.
\end{enumerate}

\subsubsection{Целевая функция (неявная оптимизация)}
\noindent
Минимизация:
\begin{itemize}
    \item Время рендеринга $T_{render}$ изображения $I$.
    \item Количество треугольников $M$ в сетке $\Delta$ при заданной визуальной точности $T_{detail}$.
\end{itemize}

\noindent
Максимизация:
\begin{itemize}
    \item Визуальное качество и точность аппроксимации $\Delta$ поверхности $S$ с учетом $T_{detail}$.
\end{itemize}

\subsubsection{Ограничения}
\begin{itemize}
    \item Реализация $\mathcal{G}$ на GPU с использованием современных кросплатформенных API (Vulkan) и аппаратных возможностей Mesh Shaders \cite{NVIDIA_MeshShaders}.
    \item Адаптивность триангуляции в $\mathcal{G}$ к видимым для $C$ участкам и уровню детализации $T_{detail}$.
\end{itemize}

\section{План решения задачи}

В рамках данной магистерской диссертации необходимо решить следующие задачи:
\begin{enumerate}
    \item 
Провести обзор существующих методов динамической триангуляции поверхностей, задаваемых функцией расстояния.

    \item 
Разработать базовый метод триангуляции в прототипе в виде программы на С++ на центральном процессоре.

    \item 
Разработать программное окружение для реализации триангуляции на графическом процессоре с использованием Vulkan и Mesh \cite{NVIDIA_MeshShaders} шейдеров.

    \item 
Разработать простейший вариант реализации алгоритма динамической триангуляции и растеризации по одной из выбранных статей.

    \item 
Разработать базовый вариант реализации алгоритма динамической триангуляции и растеризации по одной из выбранных статей. На данном этапе триангулированная поверхность может быть явно сохранена в память графического процессора.

    \item 
Разработать улучшенный вариант алгоритма динамической триангуляции и растеризации. Предполагается что на данном этап алгоритм отправляет сгенерированные треугольники непосредственно на отрисовку в графический конвейер используя возможности меш-шейдеров \cite{NVIDIA_MeshShaders}.

    \item 
Провести валидацию разработанного подхода минимум на 10 различных 3D моделях. За эталон можно взять рендеринг моделей трассировкой лучей с 4 или восьми сэмплами на пиксел (анти-алиасинг).

    \item 
Провести сравнение предложенного метода с базовой реализацией и как минимум с методом визуализации на основе трассировки лучей.

\end{enumerate}

\section{Обзор существующих методов}
Несмотря на то, что генерация SDF не является предметом диссертации, понимание их свойств остается ключевым. \textbf{Osher и Fedkiw (2003)} в своей книге \textit{Level Set Methods and Dynamic Implicit Surfaces} (\cite{Osher2003}, глава "Signed Distance Functions") по-прежнему являются фундаментальным источником для математического определения и свойств SDF. Важно понимать, что SDF, будучи функцией, где значение в точке равно кратчайшему расстоянию до поверхности (со знаком), является идеальным представлением для динамических и изменяющихся объектов. Это фундаментальная база для того, чтобы понимать, как изменение функции $\phi(x)$ приводит к изменению поверхности $\phi(x) = 0$, которую необходимо триангулировать. Мы будем работать с источником данных, который представляет собой функцию расстояния, и поэтому знание ее характеристик абсолютно необходимо.

\subsection{Методы триангуляции изоповерхностей и динамические подходы}
С учетом того, что SDF является заданным и динамически меняющимся, ключевой задачей становится эффективная экстракция постоянно обновляемой изоповерхности.

\textbf{Nesme и Bouthors (2006)} в своей диссертации \textit{Dynamic Triangulation of Implicit Surfaces: towards the handling of topology changes} \cite{nesme2006dynamic} прямо исследуют, как обновлять триангуляцию неявных поверхностей при их деформации, уделяя особое внимание обработке изменений в топологии (слияние/разделение). Это напрямую соответствует нашей теме, так как SDFы часто используются для моделирования объектов, подверженных подобным изменениям. Хотя их методы могут быть не всегда связаны с современными GPU-конвейерами, они предоставляют важные концептуальные основы для алгоритмов, отслеживающих изменения поверхности.

\textbf{Dyken et al. (2008)} в работе \textit{High-speed marching cubes using histopyramids} \cite{dyken2008high} представляют собой оптимизированную реализацию широко используемого алгоритма Marching Cubes (MC). В условиях динамической SDF, скорость MC становится критически важной. Метод гистопирамид позволяет значительно ускорить процесс нахождения ячеек, пересекающих изоповерхность, что делает его крайне релевантным для быстрого обновления триангуляции при изменении SDF. Его эффективность напрямую влияет на возможность интерактивной работы с динамическими поверхностями.

\subsection{Современные аппаратные подходы к триангуляции для динамической природы SDF}
Наиболее перспективными для динамической триангуляции, особенно в реальном времени, являются методы, активно использующие GPU.

\textbf{Elliott (2022)} в своей диссертации \textit{Using Mesh Shaders for isosurface extraction} \cite{elliott2022using} является чрезвычайно актуальным источником. Если SDF задана и меняется (например, как 3D-текстура или вычислительное поле), то Mesh Shaders (шейдеры сетки) предоставляют идеальный механизм для быстрой и гибкой генерации триангуляции \textit{полностью на GPU}. Это позволяет инкапсулировать процессы выборки значений SDF, вычисления нормалей и генерации треугольников прямо в GPU-конвейере, минимизируя узкие места, связанные с передачей данных. Эта работа показывает, как новые аппаратные возможности (описанные в блоге \textbf{NVIDIA Developer Blog (Lampe, 2020)} \cite{NVIDIA_MeshShaders}) могут быть использованы для эффективного преобразования неявного представления в явное на лету, что критически важно для динамических систем.

\textbf{Kreskowski et al. (2022)} в работе \textit{Efficient direct isosurface rasterization of scalar volumes} \cite{kreskowski2022efficient} предлагает еще более радикальный подход. Вместо того, чтобы сначала триангулировать, а затем растеризовать полученную сетку, этот метод направлен на \textit{прямую растеризацию} изоповерхности из скалярных объемов (SDFs). Это может быть наиболее эффективным подходом для \textit{визуализации} динамических SDF, поскольку полностью исключает этап построения явной геометрии, если требуется только отрисовка, а не физические взаимодействия или другие операции, требующие явной сетки. Для вашей диссертации это может представлять собой ценное сравнение или даже альтернативный метод, если конечная цель – исключительно визуализация.

\textbf{Balusek (2025)} в работе \textit{A New Interactive Algorithm For Converting a Voxel Model To a Mesh In Unreal Engine} \cite{baluseknew}, хотя и фокусируется на воксельных моделях, является показателем реальных вызовов и решений в высокопроизводительных интерактивных средах, таких как игровые движки. Принципы интерактивной конвертации "неявного" представления (как воксели, так и SDFы можно рассматривать как формы неявного представления) в явную сетку, как и оптимизации для GPU, будут крайне релевантны. Использование Unreal Engine 5 намекает на потенциал Mesh Shaders и других современных GPU-технологий, которые могут быть применены и к динамической триангуляции SDF.

\nocite{*}
\bibliographystyle{gost71u} % Для соответствия требованиям об оформлении списка литературы
\bibliography{references}

\end{document}
